\section{Backtesting Methodology}
\label{sec:backtesting}

% Lead: Rajan. Code: docs/js/cola-engine.js, docs/data/nba-data.json.

We present the first systematic historical backtest of multiple \cola{} variants using real NBA data, complementing the synthetic Bradley-Terry simulations in \citet{highley2026cola}.

\subsection{Data}

Our dataset covers 26 NBA seasons (1999--2000 through 2024--25), comprising:
\begin{itemize}
    \item 775+ team-season records from Basketball Reference, including win-loss records, playoff results (series wins/losses at each round), and draft outcomes;
    \item 776 first-round draft picks (picks 1--30, all 26 seasons);
    \item 55 pre-draft pick trades verified against lottery-day attribution (see \cref{sec:lottery-day}).
\end{itemize}

\subsection{Engine Design}

The \cola{} backtesting engine operates in two phases per season:

\begin{description}
    \item[Phase A (Pre-Draft):] Compute each team's lottery index entering the draft, based on accumulated carry-over and the current season's playoff results. This is the state displayed to users and used for draft probability calculations.
    \item[Phase B (Post-Draft):] Apply diminishment rules based on draft lottery outcomes (picks 1--4 reduce the index according to the variant's diminishment schedule). The resulting index carries forward to the next season.
\end{description}

This two-phase design resolves an off-by-one error identified by Highley in the initial implementation: without it, diminishment from the current draft would retroactively affect the probabilities used to determine that draft's outcome.

\subsubsection{Lottery-Day Attribution}
\label{sec:lottery-day}

Draft picks are frequently traded. The critical question for \cola{} index computation is: who holds the pick on \emph{lottery day} (when probabilities are determined), not on \emph{draft day} (when selections are made)?

We audited all 364 lottery picks ($14 \times 26$ seasons) and identified 8 instances where the draft-day holder differed from the lottery-day holder. The most notable: Boston held the \#1 pick on lottery day in 2017 (Philadelphia held \#3); the Fultz--Tatum swap was a post-lottery trade. Under \cola{}, Boston's index absorbs the \#1 diminishment, while Philadelphia's takes only the 50\% reduction for \#3.

\subsection{Cold Start}

\cola{} is a multi-year mechanism. Backtesting from 1999--2000 means the first several seasons have truncated drought histories. We initialize all teams with zero tickets and allow indices to build naturally. Results stabilize by approximately season 5 (2003--04). We present all 26 seasons but note the cold-start limitation.

\subsection{Variant-Specific Implementation}

\paragraph{Simple \cola{}.} 22-team eligible pool: all teams with zero playoff series wins in the current season (14 non-playoff teams + 8 first-round losers). Deterministic ordering by drought length; tiebreak by most wins.

\paragraph{Simple Lottery \cola{}.} Same drought state and 22-team pool as Simple. Pre-2019 NBA lottery odds (25.0\% down to 0.5\%) applied to the top 14 teams by drought.

\paragraph{Classic \cola{}.} Full lottery index computation per \cref{sec:cola-family}. Weighted lottery for picks 1--4; remaining picks by index rank.

\paragraph{Countdown \cola{}.} McCarty number (drought $\times$ wins) determines ranking. 5-team elimination pools with 6/5/4/3/2 tickets. Bottom-up sequential draws (21 draws total for a 22-team pool). Probabilities computed via Monte Carlo simulation per season rather than closed-form, due to the compounding effect of sequential draws. The interactive backtester runs 10{,}000 trials per season; numbers reported in case studies (Section~\ref{sec:case-mil}) come from a 100{,}000-trial reproduction with a mulberry32 PRNG seeded at 42 for paper reproducibility.

\paragraph{Capped \cola{}.} Engine implements Definition~\ref{def:capped-cola}: wins-based increment for any team finishing outside the top 6 regular-season seeds, continuous cap at $\Max$ (default 150) applied at every update step, six-step playoff diminishment (with the play-in tier separating play-in winners from play-in losers in seasons from 2019--20 onward), five-step draft diminishment, two exclusion rules (top-5 lottery winner from the prior season; playoff series winner in either the prior season or the current season), top-five lottery raffle by stockpile share, most-wins tiebreaker at the cap. The first-round playoff diminishment applies before the lottery draw because the lottery is held after round~1 (the second exclusion rule is then resolved); second-round and later diminishment applies after. To support the play-in tier, we enriched the dataset with \texttt{playInParticipant} and \texttt{playInAdvanced} fields for the 42 team-seasons in the play-in era (2019--20 bubble, 2020--21 onward), validated against the actual play-in tournament outcomes. For seasons predating the play-in tournament, the play-in tier collapses to the standard first-round-loser tier ($\rho = 0.2$).

\subsection{Limitations}

This is a \emph{static} backtest: \cola{} rules are applied to actual historical outcomes. Under a real \cola{} regime, draft outcomes would differ, which would change future drought lengths and indices. For example, the Sacramento Kings would not maintain the highest Classic \cola{} index for 6 consecutive seasons if they received the \#1 pick early in that streak. We present results with this caveat prominently noted.

Picks 15--30 reflect original holders; draft-night trades for non-lottery picks are not corrected, as these do not affect \cola{} index computation.

The Capped \cola{} backtest does not enforce two provisions of the variant's full specification: (1) the enhanced Stepien Rule (capped teams retain their own first-round pick every other year), since the pre-draft-trade dataset already contains the trade events as historical fact and re-routing them counterfactually requires the kind of dynamic backtest described in Section~\ref{sec:limitations}; and (2) the per-pick opt-out provision (a team may decline a specific pick to avoid that pick's diminishment), since the question this backtest answers is the calibration of $\Max$ rather than the equilibrium effect of strategic opt-outs. Both provisions are documented in the methodology card of the COLA Explorer~\citep{rajan2026explorer}.
